\pdfmapfile{+cm.map}
\pdfmapfile{+lm.map}
\pdfmapfile{+symbols.map}
\documentclass[11pt]{article}
\usepackage[a4paper,margin=25mm]{geometry}
\usepackage[T1]{fontenc}
\usepackage{lmodern,amsmath,amssymb,amsthm,booktabs,array,microtype,xurl}
\usepackage[hidelinks]{hyperref}
\newtheorem{theorem}{Theorem}
\newtheorem{lemma}[theorem]{Lemma}
\newcommand{\ipra}{\mathsf{iPRA}}
\newcommand{\Run}{\operatorname{Run}}
\newcommand{\Halt}{\operatorname{Halt}_{\mathrm{blank}}}
\title{Two joint 194-state designs:\\a direct relative proof in intuitionistic primitive recursive arithmetic}
\author{Research construction and verification note}
\date{14 November 2026, Pacific time}
\begin{document}
\maketitle
\begin{abstract}
We construct PA194ORMP by jointly reassociating the operand and modus-ponens
decoders of the revised PA194AX11H machine. A second construction,
PA194ORMPLC3, adds a loop-continuation reassociation, a history selector, and
three local accelerations. Both have 194 working states and preserve the
parent's complete 3,671-instruction register program. Their literal source
simulations have been rederived for all 4,778 arbitrary-register outcomes,
including arbitrary header slack and a nonempty finite spare list. We prove
directly in iPRA that each new blank-input machine has the same halting
predicate as the exact revised AX11H parent, with explicit primitive
recursive witness maps. The consistency calibration is inherited from its
specified written theorem; an end-to-end Lean formalization is not claimed.
No justified construction below 194 states was obtained.
\end{abstract}

\section{Literal objects and the exact parent}

Let $M_0$ be revised PA194AX11H, $M_1$ be PA194ORMP, and $M_2$ be
PA194ORMPLC3. Every machine has alphabet $\{0,1\}$, one two-way infinite
tape, an all-zero initial tape, and initial head position zero. Dense
working controls are $0,\ldots,193$, start is 0, and terminal control 194
has no row. Each of the 388 rows writes one bit and moves one cell left
or right. Initialization is included in the working-state count.

The corresponding native labels are 74 for start, 84 for the source
boundary, and 254 for halt. The supplied label arrays are injective
conjugacies; they add no working states. The dense tables have SHA-256:
\begin{center}\small
\begin{tabular}{p{12mm}p{135mm}}
\toprule
Machine & SHA-256\\
\midrule
$M_0$ & \path{c5dfe1026d077d3e062eca11edb789597bfbbced4f6040f8828ae5d467b1167a}\\
$M_1$ & \path{a0473681d6d5fb68e5fccfc473f35bc67aa48083d847e349ba3c52865f293629}\\
$M_2$ & \path{6d8a8cf2f5578ccdc28dd7025bc41a7c279dfbc09c5d3893cb965d70292d10f4}\\
\bottomrule
\end{tabular}
\end{center}

The parent artifacts are extracted unchanged from
\path{calibration_AX11H.zip}, whose SHA-256 is
\begin{center}\small
\path{93ee082ef035b2d32e71d6af5239b25b4bea1773f03bea4d2c39810a7fe15075}.
\end{center}
The binding checks the literal table, label map, source sites, complete
instruction array, and original contract data against that archive.
A fresh parent replay additionally agrees with every inherited affine
duration, motion, source index, branch, and symbolic-event count.

The source material is pinned to ProveIt commit
\path{7c0bb454cf73c29ca1c385f9fec856ddd236a7cb}.
The operand transformation is adapted from the incoming
\path{PA194ORH_research_2026-11-15.zip}; the MP and accelerated
loop transformations are adapted from
\path{pa194_new_designs_research_2026-11-15.zip}. Their individual
correctness results are credited to those packets. The present proof uses
fresh contracts for the joint literal tables rather than assuming that
individually valid edits automatically compose.

\section{The joint source transformations}

For the operand routine put
$A=\texttt{node32}$, $B=\texttt{node53}$,
$C=\texttt{node55}$, and $D=\texttt{node65}$.
Both new machines replace
\[
 ((A,(B,C)),D)\quad\hbox{by}\quad(A,(B,(C,D))).
\]
The resulting prefix-reader codes are
\begin{center}
\begin{tabular}{c@{\qquad}cc@{\qquad}c}
\toprule
Routine & Old word & New word & Native target\\
\midrule
$A$ & 00 & 0 & 104\\
$B$ & 010 & 10 & 107\\
$C$ & 011 & 110 & 98\\
$D$ & 1 & 111 & 90\\
\bottomrule
\end{tabular}
\end{center}
Each reader preserves every consumed tape bit. The prefix families are
complete and prefix-free.

Native control 213 also had a backward history-selector use. Omitting
that use produces an incorrect nominal 193-state table in the supplied
research. The present construction preserves it. If control 69 scans zero
with history bit $b$ immediately to its left, then
\begin{center}\small
\begin{tabular}{ccllc}
\toprule
$b$ & Steps & Parent path & New path & Common endpoint\\
\midrule
0 & 3 & $69\to213\to92\to104$ & $69\to92\to111\to104$ & $(104,h+1)$\\
1 & 2 & $69\to213\to90$ & $69\to92\to90$ & $(90,h)$\\
\bottomrule
\end{tabular}
\end{center}
Both cells are unchanged, and no exterior cell is inspected.

For the MP routine put
$E=\texttt{node75}$, $F=\texttt{node37}$,
$G=\texttt{node6}$, and $H=\texttt{node54}$.
Replace $((E,F),(G,H))$ by $(E,(F,(G,H)))$.
From control 100 the words $00,01,10,11$ are replaced respectively by
$0,10,110,111$, reaching controls $213,28,77,114$.
This reassociation contains an occurrence of the operand routine, so the
address changes of the two transformations overlap. Direct comparison
with the parent finds 154 changed addresses: ordinal intervals 224--290
and 3572--3658, inclusive. All 3,671 resolved instructions are identical.
The count is
\[
 135\text{ fixed controls}+66\text{ raw AX readers}
 -7\text{ simultaneous folds}=194.
\]

For $M_2$, the four-increment loop tail is also reassociated. Its 61
counted continuations change from depth 2 to depth 3. A caller-sensitive
124/127 selector is substituted, followed by the rows
\[
 \delta(26,1)=1R91,\qquad
 \delta(101,0)=0R17,\qquad
 \delta(190,1)=0L42.
\]
The complete final table is what the independent verifier checks.
There are 449 changed addresses relative to $M_0$. The counted-return
syntax changes at 61 sites, but every resolved target remains identical.
The argument never assumes that reassociation alone preserves counted
returns: two independent source resolvers check their actual targets.

\section{Primitive recursive configurations and complete frames}

Use a natural-number coding of finite tuples, finite lists, and signed
integers. A live source state is
\[
 z=(i,R_0,\ldots,R_8,g,Q),
\]
where $i<3671$, the registers and slack $g$ are natural numbers, and
$Q=[s]X$ is a nonempty finite list of natural numbers. A distinguished
source halt code $\dagger$ is absorbing. The source function $F$ is the
common resolved instruction array, with these auxiliary updates:
an increment appends the spare value 1 to $Q$, and a positive decrement
increases $g$ by one. Other live instructions preserve $g,Q$.
A halt instruction yields $\dagger$, and $F(\dagger)=\dagger$.
These are primitive recursive operations on valid codes. Define their
values arbitrarily, for example as $\dagger$, on malformed codes.

Let $p_j(i)$ be the finite lookup of site $i$'s address for machine $M_j$.
All three maximum address lengths are at most 37. Put
\[
 \pi=(6,2,0,4,3,7,8,5,1),\qquad V_{\pi(a)}=R_a,\qquad
 B(v)=0\,1^{v+1}.
\]
The complete frame $E_j(z)$ has control 84, head $-2$, outward left
stream $1\,0^\infty$, and right stream
\[
 p_j(i)\,0^{48+g-|p_j(i)|}\,1\,
 B(V_0)\cdots B(V_8)B(s)B(X)\,0^\infty.                 \tag{1}
\]
Set $E_j(\dagger)=\dagger$. The physical tape halves are coded by two
natural numbers $L,R$, with the low bit of $R$ scanned and the low bit of
$L$ immediately to its left. Writing $w$ and moving right sends
\[
 (L,R)\longmapsto(2L+w,\lfloor R/2\rfloor),
\]
whereas moving left sends
\[
 (L,R)\longmapsto
 (\lfloor L/2\rfloor,(L\bmod2)+2w+4\lfloor R/2\rfloor).
\]
Retain the signed head coordinate and update it by one. These equations
describe the complete tape, including the zero exteriors.

On reaching native halt 254, replace the arithmetic configuration code
by $\dagger$ and keep it fixed. This is a totalization of the arithmetic
run, not an extra physical state. Induction on literal time shows that it
has the same halting predicate as the given machine. Write
$\Run_j(t,c)$ for its primitive recursive iterate. Quantifier-free
induction proves the semigroup identity
\[
 \Run_j(u+v,c)=\Run_j(v,\Run_j(u,c))                       \tag{2}
\]
and absorption of $\dagger$.

\section{The universal physical certificates}

Each outcome duration is the scalar product of a thirteen-entry
nonnegative integer vector with
\[
 (1,g,x_0,\ldots,x_8,s,W(X)),\qquad
 W(X)=\sum_{v\in X}(v+2).                                \tag{3}
\]
For a positive decrement of register $a$, its slot represents
$x_a=R_a-1$, with $R_a=x_a+1$; other slots represent the actual
register values. Zero-decrement certificates specialize the corresponding
slot to zero. The last slot is the width of the same finite list $X$
that occurs in the frame, not an unrelated parameter.

The independent literal checker imports no arithmetic-box oracle. It
rederives four generic one-block traversals from the actual table, then
checks all 4,778 source outcomes. The observed finite data are:
\begin{center}
\begin{tabular}{lrrr}
\toprule
Quantity & $M_0$ & $M_1$ & $M_2$\\
\midrule
Source sites & 3671 & 3671 & 3671\\
Outcome certificates & 4778 & 4778 & 4778\\
Symbolic events & 1101836 & 1101819 & 1100304\\
Maximum events in one outcome & 403 & 403 & 403\\
Minimum positive constant & 3 & 3 & 3\\
Complete blank startup & 57922 & 58014 & 58010\\
\bottomrule
\end{tabular}
\end{center}

Here is the uniform arithmetic interpretation of those finite
certificates. A literal event is a substitution into the displayed
tape equations. A unary self-loop over a word of length $m$ follows by
quantifier-free induction on the processed prefix, with the untouched
contexts as parameters. A single generic block traversal has duration
$v+2$; induction on the prefix length of a fixed finite list gives
duration $W(X)$ for the concatenation. All four one-block proofs inspect
no opaque exterior context. Word normalization uses concatenation,
reversal, and the identities
\[
 C(X)1=1B(X),\qquad
 1\operatorname{rev}(C(X))=\operatorname{rev}(B(X))1,
 \qquad C(v)=10\,1^v.
\]
Their proofs are again induction on the finite list's prefix length.
The imported \path{calibration/physical_ipra_audit.tex} supplies the
full interpretation of precisely these event rules, including the
cut orientation of the left and right streams. The same rules and four
list schemas are used in the new replays. Each changed decoder row is
handled as a literal event, so no additional abstract transformation rule
is trusted.

Composition using (2) proves each successful trace uniformly in its
natural parameters and exterior list code. These induction invariants
are equalities and characteristic functions of primitive recursive
tests, with the arbitrary parameters free. No induction on a formula
with an unbounded environment quantifier is used.

The extra endpoint inspector checks every one of the 4,777 nonhalting
outcomes against both source resolvers. It verifies zero address
overflow, zero head displacement, and the literal padded-word equation
\[
 p_{\rm computed}\,0^{48-|p_{\rm computed}|}
   =p_j(i')\,0^{48-|p_j(i')|}.                             \tag{4}
\]
Appending the same $g$ slack zeros closes (1). Thus a local result for
an unnormalized address cannot be mistaken for arrival at a different
source site. The inspector checks all recorded literal and unary rows
for the absence of a transition to halt in a nonhalting outcome, and
expands each list lemma into its checked one-block rows for the same
purpose. The unique halt outcome has duration 3. The fresh parent's
row records are additionally compared with all fields of its inherited
compact contracts.

Select the relevant affine duration by the site and decidable branch,
and call it $D_j(z)$. Put $D_j(\dagger)=0$.
The certificate interpretation and (4) prove in iPRA
\[
 \Run_j(D_j(z),E_j(z))=E_j(F(z)),\qquad
 z\ne\dagger\ \Longrightarrow\ D_j(z)\ge3.               \tag{5}
\]
Source validity is preserved. The first equation remains true at
$\dagger$; the positive lower bound is deliberately restricted to live
states.

\section{Clocks, bounds, and constructive reflection}

The common initial source state is
\[
 z_*=(25,2,0,\ldots,0,0,[1]^{26}).
\]
The three independent startup checks compare the whole tape and head,
with zero exterior, and establish
\[
 \Run_j(T_j,\mathrm{blank})=E_j(z_*),\qquad
 (T_0,T_1,T_2)=(57922,58014,58010).                         \tag{6}
\]
No startup run reaches halt or leaves its checked finite tape window.

Define by primitive recursion
\[
 z_k=F^k(z_*),\qquad
 \tau_j(0)=0,\qquad
 \tau_j(k+1)=\tau_j(k)+D_j(z_k).
\]
Quantifier-free induction using (2), (5), and (6) proves
\[
 \Run_j(T_j+\tau_j(k),\mathrm{blank})=E_j(z_k).             \tag{7}
\]

For all three machines the exact coordinatewise maximum of the duration
vectors is
\[
 (534,8,4,4,4,4,4,4,2,4,4,2,2).
\]
Since each variable register slot is at most its register value, and
$s+W(X)\le W(Q)=s+2+W(X)$, this gives the safe bound
\[
 3\le D_j(z)\le534+8g+4\sum_{a<9}R_a+2W(Q)
 \quad(z\ne\dagger).                                    \tag{8}
\]
Along the run from $z_*$, primitive recursive bookkeeping and
quantifier-free induction give
\[
 g_k\le k,\qquad \sum_{a<9}R_{a,k}\le2+k,
 \qquad |Q_k|\le26+k,\qquad W(Q_k)\le78+3k.              \tag{9}
\]
Indeed, only a positive decrement increases slack, by one; each source
step increments at most one register; and only an increment appends one
spare entry, equal to 1. Initially the 26 spare blocks each have width
3. No source step removes the first spare or changes a spare value.
The width bound follows directly, and is valid independently of the
list coding. Define these selectors to be zero at $\dagger$ to extend
the inequalities beyond source halt.

Substituting (9) in (8), and using $D_j(\dagger)=0$, gives
\[
 0\le D_j(z_k)\le
 534+8k+4(2+k)+2(78+3k)=698+18k.
\]
Consequently
\[
 \tau_j(k)\le698k+18\frac{k(k-1)}2=9k^2+689k.             \tag{10}
\]
Thus 689 is $698-9$, not an empirically fitted constant. This bound
is deliberately conservative: replacing $s+W(X)$ by $W(Q)$ loses
four units in the duration bound. Put
\[
 P_j(k)=T_j+9k^2+689k.                                    \tag{11}
\]
After a halt already reached by (7), absorption gives a halt also
at the upper-bound time $P_j(k)$.

For reflection, induction on $k$ using source absorption and (5) proves
\[
 z_k\ne\dagger\quad\Longrightarrow\quad\tau_j(k)\ge3k.
                                                               \tag{12}
\]
Suppose the physical machine has halted by time $t$, and set $k=t+1$.
If $z_k\ne\dagger$, then $T_j+\tau_j(k)>t$, so physical absorption
would make the left side of (7) halted while its right side is a live
frame at control 84. This is impossible. Equality of a coded source
state with $\dagger$ is decidable, so intuitionistic reasoning gives
$z_{t+1}=\dagger$. No Markov principle or unbounded least-time search
is used. Conversely, $z_k=\dagger$ implies a physical halt by $P_j(k)$.

Write $h_j(t)$ for the primitive recursive assertion that $M_j$ is
halted by time $t$ from blank. We have proved the explicit witness
transports, for $j=1,2$,
\begin{align*}
 \ipra&\vdash h_0(t)\longrightarrow h_j(P_j(t+1)),\\
 \ipra&\vdash h_j(t)\longrightarrow h_0(P_0(t+1)).        \tag{13}
\end{align*}
The same argument even permits $\lfloor t/3\rfloor+1$ in place of
$t+1$, but that sharper bound is unnecessary for the stated theorem.

\begin{theorem}[Direct relative equivalence]
For the literal tables bound above and each $j\in\{1,2\}$,
\[
 \ipra\vdash\Halt(M_j)\leftrightarrow\Halt(M_0),
 \qquad
 \ipra\vdash\neg\Halt(M_j)\leftrightarrow\neg\Halt(M_0).
\]
\end{theorem}
\begin{proof}
Existential introduction and elimination applied to (13) give the
positive equivalence. For either negative implication, assume a halt
witness on the opposite side, transport it by (13), and contradict the
assumed nonhalting. This is intuitionistic logic. All witness functions
displayed in (11)--(13) are primitive recursive.
\end{proof}

\section{The consistency corollary and formalization boundary}

The retained revised AX11H written calibration gives
\[
 \ipra\vdash\operatorname{Con}(\mathsf{PA})
              \leftrightarrow\neg\Halt(M_0),             \tag{14}
\]
for its explicitly mapped current Foundation PA-refutation predicate.
The archive identifies the numerical calculus, the source-command
theorem, proof translation, weak-arithmetic transfer, and literal
simulation. Those are mathematical dependencies of (14), credited
to the supplied work. The present relative theorem does not silently
replace a PA-level theorem by an iPRA theorem.

\begin{theorem}[Consistency corollary using the revised calibration]
Using the exact written calibration (14), for $j=1,2$,
\[
 \boxed{\ipra\vdash\neg\Halt(M_j)
                       \leftrightarrow\operatorname{Con}(\mathsf{PA}).}
\]
\end{theorem}
\begin{proof}
Compose the negative relative equivalence with (14). This composition
does not assume the consistency of PA and uses no new arithmetic
principle.
\end{proof}

The new executable evidence checks finite bindings, whole-tape symbolic
contracts, frame closure, literal startup, and the stated clock bounds.
The written induction arguments supply their uniform iPRA meaning.
Neither component is represented as a completed Lean kernel theorem.
The explicit tables and reconstruction are unconditional artifacts;
the absolute consistency endpoint retains the identified written
calibration dependency. No bounded negative search establishes a
lower bound of 194 states.
\end{document}
